% a06_BALtorus_ct4: Step 3(E), fourth part: the vacuum counterterm pieces E_{K,k} at every step (Ba3 (1.28), (1.30)),
% Ba2 (2.14) for the constant the chain carries, and Ba1 Props. 3.1-3.2 and Ba2 p. 579 ("Prop. I.3.1 can also be applied").

\subsubsection*{The vacuum counterterm pieces}

The vacuum counterterm $E_1$ is split by \cite[(1.26), (1.28), (1.30)]{Ba3} as $E_1=\Efin+E_{K,k}+E_k$, $E_K=E[G_K,G^v_K]$, $E_k=E[G_k,G^v_k]$ (class sums \cite[(1.24)]{Ba3} with the counterterm coefficients kept as numbers, \cref{conv:a06d:LOC}), $E_{K,k}\defeq E_K-E_k$ and $\Efin\defeq E_1-E_K$. $\Efin$ is a global constant, factored out of the chain by \cref{lem:a06:CF} and bounded by \cref{thm:a06:F}; the pieces $E_{K,k}$ are what the step-$k$ action carries, and the following statements bound them with constants free of $\eps_0$.

\begin{lemma}[the $c$-number of the step-$k$ polynomial]\label{lem:a06e:VK0}
At zero background on the torus, the field-independent part of the step-$k$ polynomial \cite[(1.4)--(1.5)]{Ba3} with vertex domain $\Omega_1$ is $c_k(\Omega_1)=-(\Efin+E_{K,k})+a_k(\Omega_1)$, where $a_k(\Omega_1)$ is the sum of the connected vacuum families containing at least one block-spin pair vertex \cite[(1.12)--(1.14)]{Ba3}. On the chain run with $E_0+E_K$ (\cref{lem:a06:CF}) it is $-E_{K,k}+a_k(\Omega_1)$.
\end{lemma}

\begin{proof}
By \cite[(1.4)]{Ba3} the vacuum graphs of $-\log$ of the fluctuation integral at zero background without pair vertices have the vertices of the action localised in $\Omega_1$ (\cite[p.~419]{Ba3}: the vertices \cite[(1.6)--(1.7)]{Ba3} by $\sum_{x\in\Omega_1}$, the others by $g_k$, and ``we perform this localization for the terms of $\dm_k$, $E_k$ also'' \cite[p.~420]{Ba3}), the full $\dm$ of \cite[(1.29)]{Ba3} in the vertex \cite[(1.7)]{Ba3}, and lines $G_k(T,0)$; order by order these are the graphs of $E_k(e'e_kg_k,\lambda'\lambda_k,\Omega_1)$ (\cite[p.~418]{Ba3}: ``the same form, only the propagators are replaced by $G^\eps_k(0)$, $G^\eps_k$''), while $\Efin$ and $E_{K,k}$ carry no $\Omega_1$ \cite[(1.30), p.~419]{Ba3}. So the $\Omega_1$-dependent vacuum part cancels against $-E_k(\Omega_1)$ in \cite[(1.5)]{Ba3}, and the pair-vertex families remain.
\end{proof}

\begin{theorem}[the vacuum pieces at step $k$]\label{thm:a06e:VK}
For $0\le k<K$, on every torus aligned with the blocks with sides $\ge1$: \textup{(a)} the pair-vertex vacuum families $a_k(\Omega_1)$ are $O(t_k)$ per localisation cube; \textup{(b)} $|E_{K,k}|\le C_{\mathrm{VK}}t_k|T^{(k)}|=C_{\mathrm{VK}}\ell_T^3l_k^{-2}$ for the classes without $\dm_{(4)}$ vertex, and $\le C_\beta t_k^{1-\beta}|T^{(k)}|$ for every $\beta>0$ for the classes with one; \textup{(c)} at order $\lambda$, without Proposition~1: $E^{(0,1)}_{K,k}=\lambda N(N+2)\sum_x\eps^3G_{K,k}(x,x)[R_K(x,x)+R_k(x,x)]$, so $|E^{(0,1)}_{K,k}|\le2\lambda N(N+2)c_R^2\ell_T^3l_k^{-2}$, and $0\le\Efin^{(0,1)}=\lambda N(N+2)\sum_x\eps^3R_K(x,x)^2\le\lambda N(N+2)c_R^2\ell_T^3l_K^{-2}$. The constants are free of $\eps$, $\eps_0$, $k$, $K$ and the torus.
\end{theorem}

\begin{proof}
(a) Proposition~1 with no external fields gives $O(1)e_k^{d_v}\lambda_k^{d_s}e^{-\delta_0d}$ per localised family; a pair vertex contains a \cite[(1.14)]{Ba3} member, of order $\ge1$ in $e$, and all $A'$-legs are contracted in pairs, so $d_v\ge2$ and the family carries $t_k$ (as in the proof of \cref{lem:a06d:V9}); sum over localisations with the root fixed (\cref{lem:a06d:V11}). (b) $E_{K,k}=\sum_{j=k}^{K-1}\epsilon_j$ with $\epsilon_j\defeq E[G_{j+1}]-E[G_j]$, and \cref{lem:a06e:CUT}(d) for vacuum classes of order $o$ gives $|\epsilon^{(o)}_j|\le Ct_j^{o/2}|T^{(j)}|=C\ell_T^3l_j^{o/2-3}$; summed over $j\ge k$: $o=2$ gives $C\ell_T^3l_k^{-2}$, $o=4$ gives $C\ell_T^3l_k^{-1}$, $o=6$ gives $C\ell_T^3(K-k)$, $o\ge8$ gives $C\ell_T^3$, all $\le C\ell_T^3l_k^{-2}$ since $l_k^2(K-k)\le1$. The \cite[(1.7)]{Ba3} vertices of the vacuum families have orders $2$ or $4$; the non-renormalising part of a $\dm_1$ vertex at step $j+1$ is $\Delta_{j+1,j+1}=O(l_{j+1}^{-1})$ (\cref{prop:a06e:DM1}), of the printed size, and that of a $\dm_{(4)}$ vertex is $\sum_{i>j}m^{(4)}_i+{\dm}^{(4)}_{\mathrm{fin}}$, whose coefficient in scale-$(j+1)$ units is $\le C_Tl_{j+1}q(l_{j+1})$ (\cref{thm:a06e:T0}(e)), so such a vertex carries $t_j^{1-\beta}$ in place of $t_j^2$ and the classes containing one obey $C_\beta t_j^{o/2-1-\beta}|T^{(j)}|$; summed over $j\ge k$ with $o\ge4$ this is $\le C_\beta\ell_T^3l_k^{-2-\beta}=C_\beta t_k^{-\beta}t_k|T^{(k)}|$. (c) $E^{(0,1)}[G]=-\lambda N(N+2)\sum_x\eps^3G(x,x)[G(x,x)-2C_0(0)]$ (the figure-eight and the cycle with the $\dm_{(0,1)}=-4(N+2)C_0(0)$ vertex); subtracting at $G_K$ and $G_k$, resp.\ at $C_0$ and $G_K$, gives the two displayed identities, and \cref{lem:a06e:B0} the bounds.
\end{proof}

\begin{corollary}[{\cite[(2.14)]{Ba2} for the constant the chain carries}]\label{cor:a06e:C1b}
On the chain run with $E_0+E_K$ the vacuum constant at step $0$ is $E_K=E_{K,0}$ ($G_0=0$), and $|E_K|\le C\eps|T_\eps|$ for the classes without $\dm_{(4)}$ vertex, $\le C_\beta\eps^{1-\beta}|T_\eps|$ for the others, with constants free of $\eps$, $\eps_0$ and the torus. Read as a density, $E_K(X)\defeq(E_K/|T_\eps|)|X|$: $E_K=E_K(\Lambda_7)+E'_K(\Lambda_7^c)$ with $|E'_K(\Lambda_7^c)|\le C_\beta\eps^{1-\beta}|\Lambda_7^c|$, which is \cite[(2.14)]{Ba2} with $\kappa_0=1-\beta$ for every $\beta>0$ ($\kappa_0=1$ for the classes without $\dm_{(4)}$ vertex), for the constant to which the chain applies it.
\end{corollary}

\begin{proof}
\Cref{thm:a06e:VK}(b) at $k=0$, $t_0|T^{(0)}|=\eps|T_\eps|$; \cref{lem:a06:CF} for the booking.
\end{proof}

This bound holds for the class sums only. The vacuum graph made of a scalar ring through three $\lambda$-vertices, each with one scalar self-loop, is a term of $E_1$ of order $\lambda^3$ whose value per $\eps$-site is $\ge c_N\lambda^3m^{-3}$ for small $\eps$ (its three insertions are $4(N+2)\lambda C_0(0)\sim\eps^{-1}$ each and $\ell_T^{-3}\sum_p\hat C_0(p)^3\ge c\,m^{-3}$); it cancels within the order-$\lambda^3$ coefficient against the rings with $\dm_{(0,1)}$ insertions.

\begin{proposition}[{\cite[Props.~3.1--3.2]{Ba1} and \cite[p.~579]{Ba2}}]\label{prop:a06e:B12}
Given Proposition~1 at $L=5$ (Steps~3(D$_3$)--3(D$_4$)) with the vertex sizes of \cref{thm:a06e:T0}: \textup{(i)} (\cite[Prop.~3.1]{Ba1}) at step $k$ of the lower bound, the sum of the $R$-terms (the remainders of the Taylor expansion in the background, \cite[(1.9), (1.11), (1.15)]{Ba3}) and the sum of the field-dependent terms at whose vertices the powers of $L^k\eps$ add up to more than $d$ are $O(1)(L^k\eps)^\kappa|T^{(k)}_1|$ with $\kappa>d$; the $c$-number part of the step-$k$ expansion is $-E_{K,k}+a_k(\Omega_1)$, bounded by \cref{thm:a06e:VK}, and the total $c$-number the chain collects is $E_K=O(\eps|T_\eps|)$, which the final bound absorbs. \textup{(ii)} (\cite[Prop.~3.2, (3.57)--(3.58)]{Ba1}) the coefficients $v^{(k)}_{j;\mu}(B^{(k+1)},\psi;x,y)$ of $V^{(k)}$ with $n+m\ge1$ obey $|v^{(k)}|\le O(1)(L^k\eps)^{\kappa_0}\exp(-\delta_0d(x_1,\dots,y_m))$ with $\kappa_0>0$ independent of $k$; the coefficient with $n=m=0$ is a function of the new fields $(B^{(k+1)},\psi)$, whose field-dependent graphs obey the same bound through their external legs, and whose pure $c$-number part is that of \textup{(i)}: $O(t_k^{1-\beta})$ per localisation cube. \textup{(iii)} (\cite[p.~579]{Ba2}) at step $k$ of the upper bound, on $\Omega=B^k(\Lambda^{(k)}_0)$ with the background $\tilde B$ regular by \cite[(2.98)]{Ba2}, the terms with $R$-vertices and the terms of order $>\bar n$ are $O((L^k\eps)^\kappa)|\Lambda^{(k)}_2|$, the terms containing $\tilde B^{(k)}$ are $O(t_k^\kappa)|\Lambda^{(k)}_2|$ for every $\kappa$ by \cref{cor:a06d:G}, and the field-dependent terms of $Q^{(k)}$ localised in $B^k(\Lambda^{(k)c}_7)$ are bounded by \cref{prop:a06f:H}, the corresponding terms of $P^{(k)}$ being the subject of \cref{rem:a06f:HP}; the $c$-number part is as in \textup{(i)}.
\end{proposition}

\begin{proof}
(i) Every $R$-term carries $e(L^k\eps)^{\bar n+1}$ with $\bar n+1$ background legs, and Proposition~1 with the field restrictions of \cite[Sect.~3]{Ba1} bounds it; a field-dependent term with more than $d$ powers of $L^k\eps$ at its vertices is bounded by Proposition~1 with that power; the exponent count $\kappa>d$ is arithmetic on the printed orders and needs $\bar n\ge6$ (\cref{thm:a06:BALA}). The $c$-number part is \cref{lem:a06e:VK0} with \cref{lem:a06:CF}, bounded by \cref{thm:a06e:VK}(a),(b); the total is \cref{cor:a06e:C1b}. For a $c$-number vacuum class of order $>6$ no graph-by-graph bound of the printed kind holds (a ring of $n$ step-$k$ tadpole insertions $\lambda_kG_k(x,x)\sim\lambda t_k/\eta$ has size $\sim(\lambda t_k/\eta)^n$ per unit block), which is why the class sums of \cref{thm:a06e:VK} are used. (ii) For $n+m\ge1$ this is Proposition~1's localised bound at fixed external fields, with $\kappa_0$ its uniform exponent (\cref{lem:a06e:T1} for the \cite[(1.7)]{Ba3} vertices); the coefficient with $n=m=0$ is free of the fluctuation fields but depends on $(B^{(k+1)},\psi)$, its graphs with external legs fall under Proposition~1, and its pure $c$-number part is read per localisation cube as the density of $-E_{K,k}+a_k(\Omega_1)$. (iii) is (i)'s argument on $\Omega=B^k(\Lambda^{(k)}_0)$: Proposition~1 is stated for a general $\Omega$ with a regular background \cite[p.~412]{Ba3}, and the depth and H\"older clauses of the background are those of \cite[(2.94)--(2.98)]{Ba2}; the $\tilde B^{(k)}$ terms are \cref{cor:a06d:G}; the $Q^{(k)}$ terms are \cref{prop:a06f:H} and the $P^{(k)}$ terms \cref{rem:a06f:HP}.
\end{proof}
